% !TEX program = xelatex
\documentclass[10pt,a4paper]{ctexart}

\usepackage{amsmath,amssymb}
\usepackage{booktabs}
\usepackage{geometry}
\usepackage{enumitem}
\usepackage{longtable}
\usepackage{array}
\usepackage{xcolor}
\usepackage{xurl}
\usepackage{hyperref}

\geometry{left=1.8cm,right=1.8cm,top=2.2cm,bottom=2.2cm}
\hypersetup{colorlinks=true,linkcolor=blue,urlcolor=blue}

\newcolumntype{L}[1]{>{\raggedright\arraybackslash}p{#1}}
\urlstyle{tt}
\newcommand{\eng}[1]{\begingroup\footnotesize\path{#1}\endgroup}
% 节号 + 中文名，提高可读性
\newcommand{\modref}[2]{\textit{#1~#2}}

\title{%
  \textbf{星闪 SLB 物理层}\\[0.3em]
  \Large 6.6 模块实现参数汇总（详尽输入/输出）\\[0.2em]
  \large 角色配置 \textbullet\ RTT/FLAG/跳频 \textbullet\ AoA/AoD \textbullet\ 六路径解算
}
\author{依据 6.6 分册与接口暴露文档；airMode\,=\,\texttt{slb}\\
\small 每模块写明输出 / 直接输入 / 他模块输出引用（不展开上游入参）\\
\small 他模块一律写作``节号~名称''，避免只写裸节号
}
\date{2026 年 7 月 28 日}

\begin{document}
\maketitle
\tableofcontents
\newpage

\section{约定}
\textbf{输出}：本模块对外可引用量。\\
\textbf{直接输入}：L0/L1/L2/L3 本侧入参。\\
\textbf{他模块}：写作「节号~中文名」（必要时再跟输出工程名），例如
\modref{6.6.1}{角色与汇聚}、\modref{6.2.3.9}{PMS波形}。

%----------------------------------------------------------------------
\part{6.6.1 角色与汇聚}
%----------------------------------------------------------------------

\section{6.6.1 角色配置与测量汇聚}
\textbf{API}：\eng{slbPosCreateContext}/\eng{slbPosApplyXrcConfig}/\eng{slbPosApplyGciResource}/\eng{slbPosValidateConfig}

\begin{longtable}{L{3.2cm}L{12cm}}
\caption{6.6.1 角色与汇聚 I/O}\label{t661} \\
\toprule\textbf{角色}&\textbf{内容}\\\midrule
\endfirsthead
\toprule\textbf{角色}&\textbf{内容}\\\midrule
\endhead\bottomrule\endfoot
输出 &
\eng{SlbPosContext}（\eng{localPhyId}、能力、\eng{nodes[]} 角色/锚点坐标、\eng{frameType} ABS/REL、\eng{fixNodePhyId}、可选测量位图、静态 PMS、上报目标、组播）；
\eng{aggregatedMeasSet[]}；\eng{syncReadyFlag}；\eng{configValid} \\
直接输入 &
第7章已解析 XRC/能力（禁 raw ASN.1）；
GCI 格式0--2 解码；SAB 同步完成；对端测量 IE \\
他模块 &
\modref{第7章}{XRC/能力IE}（解析字段）；
\modref{6.2.4.6.6}{GCI位置测量资源指示}（格式0--2）；
\modref{6.2.3}{物理层信号}/SAB（同步完成指示） \\
\end{longtable}

%----------------------------------------------------------------------
\part{6.6.2 测距}
%----------------------------------------------------------------------

\section{6.6.2.1 RTT 点对点测距}
\textbf{API}：\eng{slbPosBuildRttAirPlan}/\eng{slbPosOnToaTodSample}/\eng{slbCalculateRttDistance}

\begin{longtable}{L{3.2cm}L{12cm}}
\caption{RTT I/O}\label{t6621} \\
\toprule\textbf{角色}&\textbf{内容}\\\midrule
\endfirsthead
\toprule\textbf{角色}&\textbf{内容}\\\midrule
\endhead\bottomrule\endfoot
输出 &
\eng{rttAirCommands[]}；
\eng{taUs}..\eng{tdUs}；
\eng{distanceM}；
时间差报告 IE \\
直接输入 &
\eng{rttMode} 两/三信号；
$t_1{\ldots}t_6$（Us）；
\eng{speedOfLightMps} \\
他模块 &
\modref{6.6.1}{角色与汇聚}（上下文/资源）；
\modref{6.2.3.9}{PMS波形}；
\modref{6.5.5.7}{PMS时间差TOA/TOD}；
\modref{6.2.4.6.6}{GCI位置测量资源指示}（格式1--3） \\
公式 &
两信号 $D=(T_a-T_b)C/2$；
三信号 $D=(T_a T_d-T_b T_c)/(T_a+T_b+T_c+T_d)\,C$ \\
\end{longtable}

\section{6.6.2.2 FLAG 快速定位}
\textbf{API}：\eng{slbPosSetupMeasGroup}/\eng{slbPosOnGciFormat4}/\eng{slbPosBuildFlagAirPlan}/\eng{slbPosCollectToaSet}

\begin{longtable}{L{3.2cm}L{12cm}}
\caption{FLAG I/O}\label{t6622} \\
\toprule\textbf{角色}&\textbf{内容}\\\midrule
\endfirsthead
\toprule\textbf{角色}&\textbf{内容}\\\midrule
\endhead\bottomrule\endfoot
输出 &
组上下文（组ID、PhyID序、模式 M2M/DL/SUL/AUL、信号形态、报告方向、PMS $u$、间隔）；
\eng{allowTxPms}；
\eng{flagAirCommands[]}；
\eng{distanceMatrixM} 或 \eng{tdoaInputSet[]} \\
直接输入 &
建组/参数消息；GCI 第四格式（激活+报告调度）；同步完成；多锚 TOA \\
他模块 &
\modref{6.6.1}{角色与汇聚}；
\modref{6.2.4.6.6}{GCI位置测量资源指示}（第四格式激活）；
\modref{6.2.3.9}{PMS波形}；
\modref{6.5.5.7}{PMS时间差TOA/TOD}；
可复用 \modref{6.6.2.1}{RTT测距} 的 \eng{distanceM} \\
开测 &
建组$\wedge$参数$\wedge$同步$\wedge$GCI-4 激活=1（F1/F2$\neq$开测） \\
\end{longtable}

\section{6.6.2.3 跳频复合测距 + 附录 L}
\textbf{API}：\eng{slbPosApplyHopConfig}/\eng{slbCalculateCompositeCsi}/\eng{slbResolveHopRangingDistance}

\begin{longtable}{L{3.2cm}L{12cm}}
\caption{跳频复合 I/O}\label{t6623} \\
\toprule\textbf{角色}&\textbf{内容}\\\midrule
\endfirsthead
\toprule\textbf{角色}&\textbf{内容}\\\midrule
\endhead\bottomrule\endfoot
输出 &
\eng{hopContext}（图案、$T_{\mathrm{RF}}$、CSI 模式、$P$、LBT、跳索引）；
\eng{compositeCsiH2w[]}；
拼接 CSI；
\eng{distanceM} \\
直接输入 &
配置 a)--l)；各跳双边 CSI IQ；\eng{freqHzPerHop[]}；CFO/SCO 补偿开关 \\
他模块 &
\modref{6.6.1}{角色与汇聚}；
\modref{6.2.3.9}{PMS波形}；
\modref{6.5.5.6}{PMS-CSI测量}；
\modref{8.3}{信道中心频率}；
附录 L（双向复合信道计算） \\
复合 &
$H^{2W}=H^{R}\times H^{I}$；相位仅余距离项（附录 L） \\
开测 &
配置齐套+跳频定时（\textbf{不}依赖 GCI-4） \\
\end{longtable}

%----------------------------------------------------------------------
\part{6.6.3 测角}
%----------------------------------------------------------------------

\section{6.6.3.1 AoA 与 6.6.3.2 AoD}

\begin{longtable}{L{1.8cm}L{4.8cm}L{8.6cm}}
\caption{角度测量 I/O}\label{t663} \\
\toprule\textbf{子节}&\textbf{输出}&\textbf{输入/他模块}\\\midrule
\endfirsthead
\toprule\textbf{子节}&\textbf{输出}&\textbf{输入/他模块}\\\midrule
\endhead\bottomrule\endfoot
AoA\newline(6.6.3.1) &
\eng{aoaAzimuthDeg}/\eng{aoaElevationDeg}；
\eng{aoaMeasSet[]} &
能力 \eng{aoA-Estimation}；\\
\modref{6.6.1}{角色与汇聚}；\\
\modref{6.2.3.9}{PMS波形}；\\
\modref{6.5.5.8}{AoA角度测量量} \\
AoD\newline(6.6.3.2) &
\eng{beamIndex}（$0{\ldots}N_{\mathrm{beam}}{-}1$）；
\eng{aodMeasSet[]} &
能力 \eng{aoD-Estimation}；波束表；\\
\modref{6.6.1}{角色与汇聚}；\\
\modref{6.2.3.9}{PMS波形}（须基于波束） \\
\end{longtable}

%----------------------------------------------------------------------
\part{6.6.4 解算}
%----------------------------------------------------------------------

\section{6.6.4 定位信息解算（六路径）}

\textbf{公共输出}：\eng{SlbPosFixResult}（\eng{tagXyzM[3]}、\eng{frameType}、\eng{fixPathId}）。

\begin{longtable}{L{1.4cm}L{3.6cm}L{1.6cm}L{8.6cm}}
\caption{六路径输入组合}\label{t664} \\
\toprule\textbf{路径}&\textbf{强制测量}&\textbf{要距?}&\textbf{输入来源（引用他模块输出）}\\\midrule
\endfirsthead
\toprule\textbf{路径}&\textbf{强制测量}&\textbf{要距?}&\textbf{输入来源（引用他模块输出）}\\\midrule
\endhead\bottomrule\endfoot
SA & AoA & 要 &
\modref{6.6.3.1}{AoA到达角}
+\modref{6.6.2}{距离测量}\eng{distanceM}
+锚坐标 \\
SD & AoD(\eng{beamIndex}) & 要 &
\modref{6.6.3.2}{AoD出发角}
+\modref{6.6.2}{距离测量}
+锚坐标 \\
MA & $N\times$AoA & 否 &
\modref{6.6.3.1}{AoA到达角} 测量集
+锚坐标 \\
MD & $N\times$AoD & 否 &
\modref{6.6.3.2}{AoD出发角} 测量集
+波束几何 \\
MR & $N\times D$ & 仅距 &
\modref{6.6.2}{距离测量} 多距离
+锚坐标 \\
MT & TOA 集 & 否 &
\modref{6.6.2.2}{FLAG快速定位}\eng{tdoaInputSet}
+锚坐标
+\modref{6.6.1}{角色与汇聚} 同步标志 \\
\end{longtable}

\textbf{强制}：单锚缺角或缺距 $\rightarrow$ \eng{ERR\_INCOMPLETE}；ABS/REL 与 \modref{6.6.1}{角色与汇聚} 配置一致。

\section{嵌套总览}
\begin{center}\small
第7章 XRC/能力
$\rightarrow$ \modref{6.6.1}{角色与汇聚}
$\rightarrow$ \modref{6.6.2}{距离测量} / \modref{6.6.3}{角度测量}
$\rightarrow$ \modref{6.6.1}{角色与汇聚}（汇聚）
$\rightarrow$ \modref{6.6.4}{定位信息解算}\\[0.4em]
测量原语：
\modref{6.2.3.9}{PMS波形}；
\modref{6.5.5.6}{PMS-CSI}；
\modref{6.5.5.7}{TOA/TOD}；
\modref{6.5.5.8}{AoA角度}
\end{center}

\end{document}
